\documentclass[../../../include/open-logic-section]{subfiles}

\begin{document}

\olfileid[hi]{sth}{z}{sep}
\olsection{पृथक्करण}

हम सहज समुच्चय-निर्माण के स्थान पर एक सिद्धांत से आरंभ करते हैं:

\begin{axiom}[पृथक्करण की योजना] प्रत्येक सूत्र $\phi(x)$ के लिए यह एक
अभिगृहीत है: किसी भी $A$ के लिए समुच्चय $\Setabs{x \in A}{\phi(x)}$ का
अस्तित्व है।
\end{axiom}

ध्यान दें कि यह कोई एक अभिगृहीत नहीं, बल्कि अभिगृहीतों की एक \emph{योजना}
है। पृथक्करण के \emph{अनंततः अनेक} अभिगृहीत हैं; प्रत्येक सूत्र $\phi(x)$
के लिए एक। योजना को समान रूप से (और सामान्यतः) इस प्रकार लिखा जाता है:

\begin{defish}
ऐसे किसी भी सूत्र $\phi(x)$ के लिए जिसमें ``$S$'' न हो, यह एक अभिगृहीत है:
\[
	\forall A \exists S \forall x(x \in S \liff (\phi(x) \land x \in A)).
\]
\end{defish}

\olref[sth][][]{part} के आरंभ में उल्लिखित परिपाटी के अनुसार पृथक्करण के
अभिगृहीतों में सूत्रों~$\phi$ के प्राचल हो सकते हैं।\footnote{इसका अर्थ
समझने के लिए \olref[sfr][infinite][induction]{natinductionschema} के ठीक बाद
की चर्चा देखें।}

हमारी समुच्चयों की संचयी-पुनरावर्ती अवधारणा पृथक्करण का तुरंत औचित्य देती
है। कारण देखने के लिए $A$ को कोई समुच्चय मानें। अतः $A$ किसी चरण~$S$ तक
बन चुका है (\stageshier{} से)। चूँकि $A$ चरण~$S$ पर बना, $A$ के सभी सदस्य
चरण $S$ से पहले बने थे (\stagesacc{} से)। अब विशेषतः उन सभी समुच्चयों पर
विचार करें जो $A$ के सदस्य हैं और $\phi$ को भी संतुष्ट करते हैं; स्पष्टतः
ये सभी समुच्चय भी चरण~$S$ से पहले बने थे। इसलिए चरण~$S$ पर उन्हें भी एक
समुच्चय $\Setabs{x \in A}{\phi(x)}$ में बनाया जाता है (\stagesacc{} से)।

सहज समुच्चय-निर्माण के विपरीत यह रसेल के विरोधाभास से बचता है। हम बस
$\Setabs{x}{x \notin x}$ समुच्चय के अस्तित्व का दावा नहीं कर सकते। इसके
बजाय, कोई समुच्चय~$A$ \emph{दिए जाने पर}, हम समुच्चय
$R_A = \Setabs{x \in A}{x \notin x}$ के अस्तित्व का दावा कर सकते हैं।
पर इससे केवल $R_A \notin R_A$ और $R_A \notin A$ सिद्ध होते हैं, जिनमें
से कोई बात बहुत चिंताजनक नहीं।

फिर भी पृथक्करण का एक तात्कालिक और उल्लेखनीय परिणाम है:

\begin{thm}\ollabel{thm:NoUniversalSet}
कोई \emph{सार्वत्रिक} समुच्चय नहीं है; अर्थात $\Setabs{x}{x = x}$ का
अस्तित्व नहीं है।
\end{thm}

\begin{proof}
विरोधाभास द्वारा प्रमाण के लिए मान लें कि $V$ सार्वत्रिक समुच्चय है। तब
पृथक्करण से $R = \Setabs{x \in V}{x \notin x} =
\Setabs{x}{x \notin x}$ का अस्तित्व है, जो रसेल के विरोधाभास के विरुद्ध है।
\end{proof}

सार्वत्रिक समुच्चय का अभाव---वास्तव में समुच्चयों के पदानुक्रम का खुला
अंत---संचयी-पुनरावर्ती अवधारणा के सबसे मूलभूत विचारों में से एक है। इसलिए
यह देखना उपयोगी है कि सहजतः हम एक दूसरे मार्ग से भी वहाँ पहुँच सकते थे।
सार्वत्रिक समुच्चय को स्वयं का !!a{element} होना होगा। किंतु हमारी
संचयी-पुनरावर्ती अवधारणा में प्रत्येक समुच्चय पदानुक्रम में (पहली बार) अपने
सभी !!{element}s के ठीक बाद वाले पहले चरण पर आता है। इससे निष्पन्न होता है
कि \emph{कोई भी} समुच्चय स्वयं का सदस्य नहीं है। क्योंकि स्वयं का सदस्य कोई
समुच्चय जिस चरण पर पहली बार आता, उसे उसके ठीक बाद पहली बार आना पड़ता, जो
असंगत है। (\olref[sth][spine][rank]{defnsetrank} में हम देखेंगे कि समुच्चय
की ``रैंक'' की धारणा से इस व्याख्या को अधिक कठोर कैसे बनाया जाए। किंतु ऐसा
करने से पहले कुछ और अभिगृहीत स्थापित करने होंगे।)

पृथक्करण और विस्तारिता के कुछ और परिणाम ये हैं।

\begin{prop}\ollabel{prop:emptyexists}
यदि किसी समुच्चय का अस्तित्व है, तो $\emptyset$ का अस्तित्व है।
\end{prop}

\begin{proof}
यदि $A$ समुच्चय है, तो पृथक्करण से
$\emptyset = \Setabs{x \in A}{x \neq x}$ का अस्तित्व है।
\end{proof}

\begin{prop}
किन्हीं समुच्चयों $A$ और $B$ के लिए $A \setminus B$ का अस्तित्व है।
\end{prop}

\begin{proof}
पृथक्करण से $A \setminus B = \Setabs{x \in A}{x \notin B}$ का अस्तित्व है।
\end{proof}

यह भी पता चलता है कि (लगभग) मनमाने सर्वनिष्ठों का अस्तित्व है:

\begin{prop}\ollabel{prop:intersectionsexist}
यदि $A \neq \emptyset$, तो
$\bigcap A = \Setabs{x}{(\forall y \in A)x \in y}$ का अस्तित्व है।
\end{prop}

\begin{proof}
$A \neq \emptyset$ मानें, अतः कोई $c \in A$ है। तब
$\bigcap A = \Setabs{x}{(\forall y \in A)x \in y} =
\Setabs{x \in c}{(\forall y \in A)x \in y}$, जिसका अस्तित्व पृथक्करण से है।
\end{proof}

किंतु $A \neq \emptyset$ की शर्त पर ध्यान दें; क्योंकि रिक्ततः
$\bigcap \emptyset$ सार्वत्रिक समुच्चय होता, जो
\olref{thm:NoUniversalSet} के विरुद्ध है।

\end{document}
